\begin{bankssolutions}{Solutions to Implicit Exercises for Chapter 4}

\BanksImplicitSolution{I-CH04-001}
Write a collective index $A=(\mu,x)$.  If $B_A$ were odd, its
time-ordered two-point function would be
\begin{equation*}
 G^{\mathrm F}_{AB}
 =\theta(x^0-y^0)\langle B_A B_B\rangle
 -\theta(y^0-x^0)\langle B_B B_A\rangle.
\end{equation*}
Exchanging both field labels gives
\begin{equation*}
 G^{\mathrm F}_{BA}=-G^{\mathrm F}_{AB}.
 \tag{4.I.1}
\end{equation*}
The Proca kernel has the opposite exchange property.  Its numerator is
symmetric in $\mu,\nu$, while its denominator is even in $p$.  Changing
$p\mapsto-p$ in the Fourier integral therefore gives
\begin{equation*}
 D_{\nu\mu}(y-x)=D_{\mu\nu}(x-y).
 \tag{4.I.2}
\end{equation*}
Equations (4.I.1) and (4.I.2) would force the two-point function to vanish.
This conflicts with its nonzero pole at $p^2=-\mu^2$ and the nonzero residue
on each of the $D-1$ physical polarizations (three in $D=4$).

The source argument makes the same obstruction algebraic.  A source for an
odd field is Grassmann odd.  For a symmetric kernel $D_{AB}=D_{BA}$,
\begin{align*}
 Q&=\int\dd A\,\dd B\,\eta_A D_{AB}\eta_B
   =\int\dd A\,\dd B\,\eta_B D_{BA}\eta_A\notag\\
  &=-\int\dd A\,\dd B\,\eta_A D_{AB}\eta_B=-Q,
\end{align*}
so $Q=0$.  The quadratic term that should generate the free propagator
disappears.  A real Grassmann Gaussian instead requires an antisymmetric
quadratic kernel.  The Proca differential operator and its inverse are
symmetric under exchange of their complete indices.

As a concrete check, take a rest-frame momentum $p^\mu=(\mu,\mathbf0)$.  The
spatial tensor in parentheses in the propagator is
$\eta_{ij}=\delta_{ij}$; its pole residue is proportional to $\delta_{ij}$
and has rank $D-1$ (rank three in $D=4$).
It is symmetric and nonzero, exactly as required by the $D-1$ bosonic
polarizations (three in $D=4$).  Thus the displayed propagator admits Bose quantization and
is incompatible with Fermi quantization.

\BanksImplicitSolution{I-CH04-002}
Use $\eta=\operatorname{diag}(-1,+1,\ldots,+1)$ and the Fourier convention
$\Phi(x)=\int_p\ee^{+\ii p\cdot x}\Phi(p)$, where
$\int_p=\int\dd^D p/(2\pi)^D$.  The free contraction of two Proca fields is
\begin{equation*}
 D_{\mu\nu}(p)
 =\frac{-\ii}{p^2+\mu^2-\ii\epsilon}
 \left(\eta_{\mu\nu}+\frac{p_\mu p_\nu}{\mu^2}\right).
 \tag{4.I.3}
\end{equation*}
For $\mu>0$ the Proca constraint has already been incorporated into this
inverse of the quadratic operator.  It introduces neither a gauge-fixing
vertex nor a ghost line.

Introduce sources for every free field and expand the interaction
exponential.  The $B$ part can be written
\begin{equation*}
 Z[J]=
 \frac{
  \exp\!\left[\ii\int\dd^D x\,
 \mathcal L_I\!\left(
  \frac{1}{\ii}\frac{\delta}{\delta J},
  \partial\frac{1}{\ii}\frac{\delta}{\delta J},\ldots
 \right)\right]Z_0[J]}
 {
  \left.\exp\!\left[\ii\int\dd^D x\,
 \mathcal L_I\!\left(
  \frac{1}{\ii}\frac{\delta}{\delta J},
  \partial\frac{1}{\ii}\frac{\delta}{\delta J},\ldots
 \right)\right]Z_0[J]\right|_{J=0}}.
 \tag{4.I.4}
\end{equation*}
Fermion sources require the corresponding ordered left and right
derivatives.  Wick's theorem pairs the derivatives in (4.I.4), producing
the following rules.

For an internal $B$ line of momentum $p$, insert (4.I.3).  Other internal
lines carry their free propagators.  For a local interaction vertex with
incoming fields $\Phi_{a_1}(p_1),\ldots,\Phi_{a_n}(p_n)$, insert
\begin{equation*}
 \mathcal V_{a_1\ldots a_n}(p_1,\ldots,p_n)
 =\left.
 \ii\frac{\delta^n S_I}
 {\delta\Phi_{a_1}(p_1)\cdots\delta\Phi_{a_n}(p_n)}
 \right|_{\Phi=0},
 \tag{4.I.5}
\end{equation*}
after stripping off $(2\pi)^D\delta^D(\sum_rp_r)$.  A derivative acting on
an incoming field supplies $+\ii p_\rho$ with the stated Fourier convention.
Several identical fields cause (4.I.5) to symmetrize their indices and
momenta automatically.

Momentum is conserved at each vertex.  Integrate $\dd^D\ell/(2\pi)^D$ for
each independent loop momentum.  A closed fermion loop supplies a minus
sign, as does every odd permutation needed to restore the selected ordering
of external fermions.  Multiply a graph by
\begin{equation*}
 S_G=\frac{1}{|\operatorname{Aut}G|},
 \tag{4.I.6}
\end{equation*}
where the automorphisms preserve field species, derivative labels, line
orientation, and all labeled external legs.  This is the residual factor
after the $1/V!$ factors in the exponential have been combined with the Wick
contractions.

Differentiating $Z$ retains every graph.  The functional
$W=-\ii\log Z$ retains the connected graphs.  The Legendre transform
$\Gamma$ retains the one-particle-irreducible graphs.  A connected Green
function has one free propagator between each source point and its first
vertex.  Amputation removes those external propagators.

For the test interaction
\begin{equation*}
 \mathcal L_I=-\frac{\lambda}{n!}
 C^{\mu_1\ldots\mu_n}B_{\mu_1}\cdots B_{\mu_n},
\end{equation*}
functional differentiation gives
\begin{equation*}
 \mathcal V^{\nu_1\ldots\nu_n}
 =-\ii\lambda C^{(\nu_1\ldots\nu_n)}.
 \tag{4.I.7}
\end{equation*}
When $C$ is symmetric, the $n!$ contractions cancel the displayed $1/n!$,
leaving $-\ii\lambda C^{\nu_1\ldots\nu_n}$.  At first order, joining all
$n$ legs of (4.I.7) to source points reproduces one external factor (4.I.3)
on every leg.  Setting $\lambda=0$ leaves the free Proca propagator, which
checks both the expansion and its normalization.

\BanksImplicitSolution{I-CH04-003}
We use the conventions of the preceding solution.  It is convenient to
record the three ingredients once:
\begin{align*}
 \Delta(r)&=\frac{-\ii}{r^2+m^2-\ii\epsilon},\notag\\
 D_{\mu\nu}(r)&=
 \frac{-\ii}{r^2+\mu^2-\ii\epsilon}
 \left(\eta_{\mu\nu}+\frac{r_\mu r_\nu}{\mu^2}\right),\notag\\
 V^\mu(p,q)&=-\ii g(p-q)^\mu,
 \qquad k+p+q=0.
 \tag{4.I.8}
\end{align*}
Here $m$ is the scalar mass, $p$ is the incoming momentum on a $\phi$
half-line, $q$ is the incoming momentum on a $\phi^*$ half-line, and $k$
enters on the vector half-line.
The last rule follows directly from $\partial\mapsto+\ii p$:
\begin{equation*}
 \ii\left[\ii g\bigl(+\ii p^\mu-\ii q^\mu\bigr)\right]
 =-\ii g(p-q)^\mu.
\end{equation*}
Every scalar propagator is oriented from a $\phi$ insertion to a
$\phi^*$ insertion.

Before listing graphs, a counting identity gives a completeness check.  Let
$b$ be the number of external $B$ legs and let the charged external legs
consist of $n$ copies of $\phi$ and $n$ copies of $\phi^*$.  For a connected
graph with $V$ cubic vertices,
\begin{equation*}
 I_B=\frac{V-b}{2},
 \qquad I_\phi=V-n,
 \qquad L=I_B+I_\phi-V+1.
\end{equation*}
Consequently,
\begin{equation*}
 V=b+2n+2L-2=E+2L-2.
 \tag{4.I.9}
\end{equation*}
Trees have $V=E-2$, and one-loop graphs have $V=E$.  Nonnegative integral
$I_B,I_\phi$ removes field assignments that cannot occur.

We first give compact drawings of the full census.  The symbol
$[B_{\sigma(1)}\cdots B_{\sigma(r)}]_\phi$ denotes an oriented scalar
$r$-gon whose external vector legs occur in the shown cyclic order.  The
symbol
\begin{equation*}
 \langle\phi(p)|B_{\sigma(1)}\cdots
 B_{\sigma(s)}|\phi^*(q)\rangle_\phi
\end{equation*}
denotes an oriented scalar line with the displayed vector insertions in
order.  A label $\Sigma_X@e$ draws a self-energy bubble on the specified
line $e$, while $\mathcal T@v$ draws the triangle correction in place of a
cubic vertex $v$.  This incidence notation fixes every edge and is
unambiguous for the labeled external legs.

\begin{center}
\small
\begin{tabular}{@{}p{.14\linewidth}p{.35\linewidth}p{.42\linewidth}@{}}
external legs & tree level & one-loop 1PI cores \\ \hline
$BB$ & free $D_{\mu\nu}$
 & $[B_1B_2]_\phi$ \\
$\phi\phi^*$ & free $\Delta$
 & one two-edge $(B,\phi)$ bubble \\
$BBB$ & none
 & $[B_1B_2B_3]_\phi$, $[B_1B_3B_2]_\phi$ \\
$B\phi\phi^*$ & one $V^\mu$
 & one triangle with internal edges $(B,\phi,\phi)$ \\
$BBBB$ & none
 & $[B_{\sigma(1)}B_{\sigma(2)}B_{\sigma(3)}B_{\sigma(4)}]_\phi$,
   $6$ cyclic orders \\
$BB\phi\phi^*$
 & $\langle\phi|B_1B_2|\phi^*\rangle_\phi$ and
   $\langle\phi|B_2B_1|\phi^*\rangle_\phi$
 & two boxes, according to the order of $B_1,B_2$ on the scalar path \\
$\phi\phi\phi^*\phi^*$
 & two charge-respecting vector exchanges
 & four alternating $(B,\phi,B,\phi)$ boxes
\end{tabular}
\end{center}

Every graph in this table and every labeled insertion described below has
symmetry factor one.  With external labels fixed, the vertex permutations
cancel the factors from the interaction exponential.  Scalar orientation
distinguishes the two lines in the vector self-energy bubble.  For a pure
vector polygon, a cyclic order has already been quotiented by rotations;
the reversed order is a separate oriented contraction.

Here are momentum routings and the unevaluated integrals for the loop cores.
The two self-energies are
\begin{align*}
 \mathcal B_{\mu\nu}(k)
 &=\int_\ell
 V_\mu(\ell,-\ell-k)\Delta(\ell+k)
 V_\nu(\ell+k,-\ell)\Delta(\ell),\notag\\
 \mathcal S(p)
 &=\int_\ell
 V^\alpha(p,-p-\ell)D_{\alpha\beta}(\ell)
 \Delta(p+\ell)V^\beta(p+\ell,-p),
 \tag{4.I.10}
\end{align*}
where $\int_\ell=\int\dd^D\ell/(2\pi)^D$.  The connected two-point
corrections, with their external lines restored, are
$D\mathcal B D$ and $\Delta\mathcal S\Delta$.

For the one-loop $B\phi\phi^*$ vertex, take all external momenta incoming,
$k+p+q=0$.  A useful triangle routing is
\begin{align*}
 \mathcal T^\mu(k;p,q)
 =\int_\ell&
 V^\alpha(p,-p-\ell)D_{\alpha\beta}(\ell)
 \Delta(p+\ell)\notag\\
 &\times V^\mu(p+\ell,q-\ell)
 \Delta(\ell-q)V^\beta(\ell-q,q).
 \tag{4.I.11}
\end{align*}
The three-point connected function contains the tree vertex, this triangle,
and one external self-energy insertion on each of its $B$, $\phi$, and
$\phi^*$ legs.  Thus it has one tree graph and four nonzero one-loop graphs.
Each external insertion means multiplication by the appropriate propagators
and by $\mathcal B$ or $\mathcal S$ from (4.I.10).

For pure vector loops, choose a cyclic order $\sigma$ and put
\begin{equation*}
 r_1=\ell,
 \qquad r_{j+1}=r_j+k_{\sigma(j)},
 \qquad r_{r+1}=r_1.
\end{equation*}
The oriented scalar $r$-gon is
\begin{equation*}
 \mathcal C^{(\sigma)}{}_{\mu_1\ldots\mu_r}
 =\int_\ell\prod_{j=1}^r
 \left[
   \Delta(r_j)\,
   -\ii g(r_j+r_{j+1})_{\mu_{\sigma(j)}}
 \right].
 \tag{4.I.12}
\end{equation*}
For $r=3$, the two cyclic orders are $(123)$ and $(132)$.  Reversing the
charge-flow orientation and shifting $\ell$ changes the sign of each of the
three current vertices.  Hence
\begin{equation*}
 \mathcal C^{(123)}{}_{\mu_1\mu_2\mu_3}
 +\mathcal C^{(132)}{}_{\mu_1\mu_2\mu_3}=0.
 \tag{4.I.13}
\end{equation*}
The one-loop $BBB$ function vanishes by charge conjugation.  For $r=4$,
there are $(4-1)!=6$ cyclic orders.  They form three reversal pairs and add
rather than cancel.  These six boxes are the complete connected $BBBB$
answer through one loop.

The tree-level $BB\phi\phi^*$ function also displays its two drawings
directly.  Let $k_1,k_2,p,q$ be incoming, with $p$ on $\phi$, $q$ on
$\phi^*$, and $k_1+k_2+p+q=0$.  Its amputated value is
\begin{align*}
 \mathcal A^{\mathrm{tree}}{}_{\mu_1\mu_2}
 =\sum_{\sigma\in S_2}&
 -\ii g(2p+k_{\sigma(1)})_{\mu_{\sigma(1)}}
 \Delta(p+k_{\sigma(1)})\notag\\
 &\times-\ii g\bigl(2p+2k_{\sigma(1)}
 +k_{\sigma(2)}\bigr)_{\mu_{\sigma(2)}}.
 \tag{4.I.14}
\end{align*}
The two one-loop 1PI boxes preserve this scalar path and join its endpoints
by an internal vector.  Give that vector incoming momentum $\ell$ at the
$\phi$ endpoint.  For either $\sigma\in S_2$ the integral is
\begin{align*}
 \mathcal X^{(\sigma)}{}_{\mu_1\mu_2}
 =\int_\ell&
 V^\alpha(p,-p-\ell)D_{\alpha\beta}(\ell)
 \Delta(p+\ell)\notag\\
 &\times -\ii g\bigl(2(p+\ell)+k_{\sigma(1)}\bigr)_{\mu_{\sigma(1)}}
 \Delta(p+\ell+k_{\sigma(1)})\notag\\
 &\times -\ii g\bigl(2(p+\ell+k_{\sigma(1)})
 +k_{\sigma(2)}\bigr)_{\mu_{\sigma(2)}}
 \Delta(\ell-q)V^\beta(\ell-q,q).
 \tag{4.I.15}
\end{align*}

Each of the two trees in (4.I.14) produces seven one-loop reducible graphs:
insert $\mathcal B$ on either external vector leg, insert $\mathcal S$ on
either external scalar leg, insert $\mathcal S$ on the internal scalar line,
or replace either cubic vertex by (4.I.11).  Together with the two boxes
(4.I.15), this gives $2\cdot7+2=16$ nonzero labeled one-loop
$BB\phi\phi^*$ diagrams.  For example, the internal-line correction is the
corresponding term of (4.I.14) with
\begin{equation*}
 \Delta(r)\longrightarrow
 \Delta(r)\mathcal S(r)\Delta(r).
 \tag{4.I.16}
\end{equation*}
The external and vertex replacements work in the same way, so their loop
momenta are already fixed by (4.I.10) and (4.I.11).

For the four-scalar function, label the incoming fields
$\phi(p_1),\phi(p_2),\phi^*(q_1),\phi^*(q_2)$.  The two vector-exchange trees
are
\begin{equation*}
 \mathcal A^{\mathrm{tree}}{}_{4\phi}
 =\sum_{\sigma\in S_2}
 V^\alpha(p_1,q_{\sigma(1)})
 D_{\alpha\beta}(p_1+q_{\sigma(1)})
 V^\beta(p_2,q_{\sigma(2)}).
 \tag{4.I.17}
\end{equation*}
The four one-loop 1PI boxes can be drawn without ambiguity by naming the
vertices $P_i,Q_j$ after their external legs.  Choose either scalar matching
\begin{equation*}
 M_\phi^{(\sigma)}
 =\{P_1Q_{\sigma(1)},P_2Q_{\sigma(2)}\},
 \qquad \sigma\in S_2.
\end{equation*}
The internal vector lines form either of the other two perfect matchings of
$\{P_1,P_2,Q_1,Q_2\}$.  The union is an alternating four-cycle.  Two scalar
matchings and two allowed vector matchings give four labeled boxes.

For example, take scalar edges $P_1Q_1,P_2Q_2$ and vector edges
$P_1P_2,Q_1Q_2$.  With momentum $\ell$ entering $P_1$ on the first vector
edge and $r=-p_1-q_1-\ell$ on the second, its integral is
\begin{align*}
 \mathcal Y={}&\int_\ell
 V^\alpha(p_1,-p_1-\ell)D_{\alpha\beta}(\ell)
 V^\beta(p_2,\ell-p_2)\notag\\
 &\times\Delta(p_1+\ell)\Delta(p_2-\ell)
 V^\gamma(p_1+\ell,q_1)D_{\gamma\delta}(r)
 V^\delta(p_2-\ell,q_2).
 \tag{4.I.18}
\end{align*}
The other three integrals follow by the stated matchings, with momentum
conservation at each vertex.  This prescription supplies their routing
without a second independent loop momentum.

Each tree in (4.I.17) again has seven one-loop reducible corrections.  Four
place $\mathcal S$ on an external scalar leg, one places $\mathcal B$ on the
internal vector, and two replace an endpoint by $\mathcal T$.  Adding the
four boxes gives $2\cdot7+4=18$ nonzero labeled one-loop
$\phi\phi\phi^*\phi^*$ diagrams.

Self-contraction of the scalar pair at one cubic vertex makes a vector
tadpole.  Momentum conservation sets its external vector momentum to zero,
and its factor is
\begin{equation*}
 T^\mu=\int_\ell V^\mu(\ell,-\ell)\Delta(\ell)
 =-2\ii g\int_\ell\ell^\mu\Delta(\ell)=0
 \tag{4.I.19}
\end{equation*}
for a translation- and Lorentz-invariant regulator.  Normal ordering gives
the same result immediately.  For completeness, the following is the census
of all connected, labeled one-loop contractions before zero subgraphs and
charge-conjugate pairs are removed:

\begin{center}
\small
\begin{tabular}{@{}lrrl@{}}
external legs & total & surviving & zero or cancelled contributions \\ \hline
$BB$ & $1$ & $1$ & none \\
$\phi\phi^*$ & $2$ & $1$ & $1$ tadpole attachment \\
$BBB$ & $2$ & $0$ & the reversal pair in (4.I.13) \\
$B\phi\phi^*$ & $6$ & $4$ & $2$ tadpole attachments \\
$BBBB$ & $6$ & $6$ & none \\
$BB\phi\phi^*$ & $24$ & $16$ & $6$ tadpoles and $2$ odd-loop attachments \\
$\phi\phi\phi^*\phi^*$ & $26$ & $18$ & $8$ tadpole attachments
\end{tabular}
\end{center}

Here is an incidence description of every tadpole entry in the last column.
Call the self-contracted vertex $T$ and join its remaining $B$ half-edge to a
vertex $U$.  For $\phi\phi^*$, both external scalar legs end on $U$.  For
$B\phi\phi^*$, a scalar edge joins $U$ to the external-$B$ vertex, and the
two orientations place the external $\phi$ and $\phi^*$ on the two unused
scalar half-edges.  For $BB\phi\phi^*$, the two remaining scalar edges form
a spanning tree on $U$ and the vertices carrying $B_1,B_2$; filling the two
unused scalar half-edges with $\phi,\phi^*$ gives the six labeled cases.  For
$\phi\phi\phi^*\phi^*$, the remaining vector edge and scalar edge join the
component containing $T,U$ to the other two vertices; the charge-respecting
placements of the four labeled external scalars give eight cases.  Each has
symmetry factor one.  Its loop momentum and zero-momentum bridge are fixed by
(4.I.19), while momentum conservation fixes every edge outside the tadpole.

There are two further vanishing $BB\phi\phi^*$ diagrams.  Put the external
$\phi,\phi^*$ on one cubic vertex and join its vector half-edge to the third
vector leg of an oriented scalar triangle; the other two vector legs carry
$B_1,B_2$.  If $K=p+q=-(k_1+k_2)$, their sum is proportional to
\begin{equation*}
 V^\rho(p,q)D_{\rho\sigma}(K)
  \left[\mathcal C^{(123)}{}_{\mu_1\mu_2}{}^\sigma
       +\mathcal C^{(132)}{}_{\mu_1\mu_2}{}^\sigma\right]_{k_3=K}=0,
 \tag{4.I.20}
\end{equation*}
by the same charge-conjugation reversal as (4.I.13).  Both diagrams have
symmetry factor one, and their loop routing is the $r=3$ routing in
(4.I.12).  Thus every connected contraction occurs either in the nonzero
list or in this vanishing census.  Equation (4.I.9), together with the
incidence patterns above, proves that the catalogue is exhaustive.

The full four-point correlators contain disconnected propagation as well.
At leading order these pieces are the three pairings $DD$ for $BBBB$, the
product $D\Delta$ for $BB\phi\phi^*$, and the two charge-respecting products
$\Delta\Delta$ for $\phi\phi\phi^*\phi^*$.  Through one loop, replace one
factor at a time by its correction in (4.I.10).  Vacuum bubbles cancel
against the normalization by $Z[0]$.  In the limit $g\to0$, the interaction
graphs disappear and these free Wick contractions remain, which checks the
census.

\BanksImplicitSolution{I-CH04-004}
The combination $B_\mu=A_\mu-\partial_\mu\theta$ is invariant under the
stated simultaneous transformation.  The source term varies by
\begin{align*}
 \delta_\Omega S_J
 &=\int\dd^D x\,J^\mu\partial_\mu\Omega\notag\\
 &=\int\dd^D x\,\partial_\mu(J^\mu\Omega)
   -\int\dd^D x\,\Omega\,\partial_\mu J^\mu.
 \tag{4.I.21}
\end{align*}
The first integral is the boundary flux
$\int_{\partial M}\dd\Sigma_\mu J^\mu\Omega$.  It vanishes when the gauge
parameter has compact support.  Gauge invariance for every such $\Omega$
then gives
\begin{equation*}
 \int\dd^D x\,\Omega(x)\partial_\mu J^\mu(x)=0.
\end{equation*}
The fundamental lemma for test functions implies
$\partial_\mu J^\mu=0$ as a distribution.  Conversely, current conservation
sets the volume term in (4.I.21) to zero.  Sufficient falloff of either
$J^\mu$ or $\Omega$ removes the boundary flux, proving invariance.

Since $A_\mu=B_\mu+\partial_\mu\theta$, the same integration by parts gives
\begin{align*}
 \int\dd^D x\,J^\mu A_\mu
 &=\int\dd^D x\,J^\mu B_\mu
   +\int\dd^D x\,J^\mu\partial_\mu\theta\notag\\
 &=\int\dd^D x\,J^\mu B_\mu
   -\int\dd^D x\,\theta\partial_\mu J^\mu
   +\int_{\partial M}\dd\Sigma_\mu J^\mu\theta.
\end{align*}
Current conservation and the same falloff condition yield
\begin{equation*}
 \int J^\mu A_\mu=\int J^\mu B_\mu.
 \tag{4.I.22}
\end{equation*}

In momentum space, conservation reads $p_\mu J^\mu(p)=0$.  Contracting two
sources with the Proca propagator produces
\begin{equation*}
 J^\mu(-p)D_{\mu\nu}(p)J^\nu(p)
 =\frac{-\ii}{p^2+\mu^2-\ii\epsilon}
 \left[J_\mu(-p)J^\mu(p)
 +\frac{(p\cdot J(-p))(p\cdot J(p))}{\mu^2}\right].
\end{equation*}
The longitudinal term vanishes.  The remaining expression has the smooth
limit
\begin{equation*}
 \lim_{\mu\to0}J^\mu(-p)D_{\mu\nu}(p)J^\nu(p)
 =\frac{-\ii J_\mu(-p)J^\mu(p)}{p^2-\ii\epsilon},
\end{equation*}
which is the conserved-source Maxwell kernel.  This check ties the
integration-by-parts condition to the finite massless limit discussed in the
surrounding text.

\end{bankssolutions}
